\begin{lemma}[Proved: a Bockstein for a linear nilpotent direction]
\label{C:linear_nil_Bockstein}
Let $A$ be a characteristic-zero DVR with maximal ideal $(2)$,
write $F=A/2A$, and let $J$ be an $A$-module on which multiplication
by $2$ is injective. Give $k=A\ltimes J$ the square-zero multiplication
$(a,x)(a',x')=(aa',ax'+a'x)$.
For each $F$-linear map $\ell:J/2J\to F$ there are natural odd maps
$b_\ell^i$ on the two Tate groups of $M/JM$, for every finite free
$k[C_2]$-representation $M$. If $u=b_\ell^2$, then for every monic
irreducible $q\in F[u]$ other than $u$ the function
\[
 \chi_{q,\ell}(M)=
 \operatorname{length}_{F[u]}
       \widehat H^0(C_2,M/JM)_q\pmod2                 \tag{C.1.25}
\]
is additive on admissible short exact sequences.
No finite-generation hypothesis is imposed on $J$.
\end{lemma}
\begin{proof}
Since every nonzero element of $A$ is a unit times a power of $2$,
the hypothesis makes $J$ torsion-free and hence flat over the DVR.
For a finite free $A$-lattice $L$ with involution the natural maps
\[
 \widehat H^i(C_2,L)\otimes_F(J/2J)
       \longrightarrow \widehat H^i(C_2,L\otimes_AJ)   \tag{C.1.26}
\]
are isomorphisms. Flatness permits passage of homology through
tensoring, and the Tate groups of $L$ are killed by $2$.
Equivalently, the plus, minus and regular decomposition (C.1.14)
checks the assertion on each of its three summands. This also proves
that the comparison maps are the natural ones, regardless of the
decomposition chosen.

Put $L=M/JM$. The local ring $k$ has maximal ideal $(2)\ltimes J$,
so its finite projective modules are free. Choose a $k$-linear
identification $M\simeq k\otimes_A L$ inducing the identity modulo
$J$, and write its involution as $T=T_0+U$. Here $T_0$ is the
involution of $L$ and
$U:L\to L\otimes_AJ$ is $A$-linear. The equations are
$T_0^2=1$ and $T_0U+UT_0=0$.
The map $U$ therefore induces an odd map between the Tate groups
of $L$ and $L\otimes_AJ$. Compose it with the inverse of (C.1.26)
and $1\otimes\ell$; this defines $b_\ell^i$.

For independence, another identification congruent to the identity
is $1+X$, where $X$ has entries in $J$. It replaces $U$ by
$U+T_0X-XT_0$. On a plus cycle the extra term is a
$(T_0-1)$-boundary, and on a minus cycle it is a
$(T_0+1)$-boundary. Thus the induced odd map is unchanged.
For a morphism $f:M\to N$, write $f=f_0+f_1$ using the split
coefficient ring. Its constant coefficient intertwines $T_{0,M}$
and $T_{0,N}$, and its $J$-coefficient obeys
\[
 U_Nf_0-f_0U_M=f_1T_{0,M}-T_{0,N}f_1.
\]
The right side is again a Tate boundary on each cycle. Consequently
$b_\ell$ is natural; this calculation also covers arbitrary changes
of the chosen basis, whose constant part need not be the identity.

An admissible short exact sequence is split over $k$ and reduces to
an exact sequence of free $A$-lattices. In an underlying splitting,
write the off-diagonal action block as $X_0+X_J$. Its $J$-coefficient
in the two-periodic differential identity is
\[
 U_AX_0+d_{A,0}^{i+1}X_J+X_Jd_{B,0}^i+X_0U_B=0.
\]
On a cycle the middle terms give a boundary in the Tate complex
with coefficients in $J$. Apply (C.1.26) and $1\otimes\ell$.
This proves the connecting-map relation
$b_{\ell,A}^{i+1}\partial^i=-\partial^{i+1}b_{\ell,B}^i$.
Thus $u=b_\ell^2$ commutes with the full exact Tate six-cycle.
On a $q$-primary part with $q\ne u$, both odd Bocksteins are
isomorphisms. The exact three-cycle and image-length argument
of Lemma~\ref{C:epsilon_Bockstein} now applies verbatim and gives
(C.1.25). All Tate spaces before tensoring with $J$ are finite
dimensional; no dimension or finite-generation assumption on $J$
has entered the argument.
\end{proof}

\begin{theorem}[Proved: an intrinsic criterion for all dyadic square-zero extensions]
\label{C:squarezero_intrinsic_criterion}
Let $A$ be any characteristic-zero DVR with maximal ideal $(2)$,
and let $J$ be an arbitrary $A$-module. For the commutative local
ring $k=A\ltimes J$ the following are equivalent:
\begin{enumerate}
\item The induction ring $\bigoplus_{n\ge0}R_k(S_n)$ is the
polynomial ring over $\mathbb Z$ on the specified $\rho_i$.
\item Its degree-two polynomial map is surjective.
\item Multiplication by $2$ is bijective on $J$.
\end{enumerate}
If these conditions fail, already the degree-two generic-comparison
kernel contains an explicit nonzero class. In particular, this
criterion applies to every square-zero thickening specified by a
surjection $k\twoheadrightarrow\mathbb Z_{(2)}$ with square-zero
kernel, without a noetherian or finite-generation assumption.
\end{theorem}
\begin{proof}
The implication (1)$\Rightarrow$(2) is immediate. If (3) holds,
every odd integer is already invertible on $J$, since it is a unit
of $A$. Hence every positive integer acts bijectively on $J$.
Proposition~\ref{C:divisible_nil_invariance} identifies all the
representation groups, induction products and specified generators
over $k$ with those over $A$. The DVR result C:DVR then proves (1).

Suppose (3) fails. First assume that $0\ne x\in J$ and $2x=0$.
The free rank-one representation $L_x$ with involution $1+x$ is
well defined, since $(1+x)^2=1$. Its determinant is $1+x$.
The determinant of an involution defines an additive homomorphism
from $R_k(C_2)$ to the multiplicative abelian group
$\{a\in k^\times:a^2=1\}$: in an admissible sequence choose an
underlying splitting and take the determinant of its triangular
action matrix. The degree-two polynomial subgroup is generated
by the constant trivial and sign lines, whose determinants lie in
$\{1,-1\}$. Since $x\ne0$ and $1+x$ has constant coefficient $1$,
it is neither $1$ nor $-1$. Thus $[L_x]$ is outside that subgroup.
Moreover $[L_x]-1$ vanishes after inverting $2$, because then $x=0$.

It remains to treat the case where multiplication by $2$ is
injective but not surjective. Choose $x\in J$ with nonzero image
in $J/2J$, and choose an $F$-linear functional
$\ell:J/2J\to F$ taking that image to $1$. Define
\[
 M_x=\left(k^2,
       \begin{pmatrix}1&x\\x&-1\end{pmatrix}\right),
 \qquad \eta_x=[M_x]-1-s.                            \tag{C.1.27}
\]
The displayed matrix squares to the identity because $J^2=0$.
Its reduction over $A$ is the direct sum of the plus and minus
lines; both maps $b_\ell$ between their one-dimensional Tate
groups are the identity. Hence
$\chi_{u-1,\ell}(\eta_x)=1$, whereas this invariant vanishes on
both constant lines. This proves that (2) fails. In fact the same
sign-conjugating matrix and evaluation sequence as for (C.1.11)
show that $\eta_x$ has exact order two.
After inverting $2$, both eigenspaces of $M_x$ are free of rank one
over the local square-zero ring $\operatorname{Frac}(A)\ltimes
(J\otimes_A\operatorname{Frac}(A))$; thus $\eta_x$ is generically
zero. This completes the equivalence and the generic-kernel assertion.

Finally, if $k\twoheadrightarrow\mathbb Z_{(2)}$ has square-zero
kernel $J$, every odd integer is a unit in $k$ because it is a unit
modulo the nilpotent ideal. The canonical map from $\mathbb Z$
therefore extends uniquely to a section
$\mathbb Z_{(2)}\to k$. The ring is consequently the split extension
$\mathbb Z_{(2)}\ltimes J$ considered above.
\end{proof}
